% Transcription — fonds Alexandre Grothendieck, shelfmark 100, batch 1 (pages 1-20).
% Established from the facsimile archives/batches/100/batch-01.pdf (a mirror of
% https://grothendieck.umontpellier.fr/100.pdf), per the skill
% transcribe-grothendieck. The batch file carries no cover sheet: PDF page N
% is page N of the folder (the Montpellier footer prints « 1 » ... « 20 »),
% and the archivists' pencil numbers bottom-left agree wherever legible
% (« 3 », « 4 », « 6 », « 8 », « 10 », « 12 », « 14 », « 16 », « 18 », « 20 »).
%
% Pass: Opus 5.5 (claude-opus-5-5), 2026-09-23 — first pass, unchecked against the pages by a human.
%
% The folder sits in the inventory group "69-89" « Géométrie et topologie
% combinatoire », dated 1976-[vers 1986], whose subject does not match this
% folder's; the folder has its own date, which \dating{} carries verbatim.
% Nothing is concluded here from the group.
%
% Fifteen of the twenty pages are transcribed. Absent:
%   - page 1, a sheet of buff card blank but for his words top right,
%     « Bonne réduction et gpes p-divisibles » — a cover; the words head the
%     \section (hand.md §5);
%   - pages 11, 13, 17 and 19, versos with nothing of his on them that bears
%     on these notes: 11 is a typed English errata list for some other text
%     (« Errata, Page 2 » ... « Page 75 »), 13 is page « -0.19- » of the
%     English typescript described below, carrying no mark of his hand
%     (checked at 150 dpi, turned upright); 17 is a typed errata list for
%     « EGA IV 4e (2èmes épreuves) », pp. 164-271, and 19 the printed cover of
%     Publications Mathématiques n° 32 with a « Bon à tirer » stamp, a date
%     stamp « 14 MARS 1967 » and circled section numbers (17.1.6, 19.1.8,
%     19.4.11, 18.12.1). They are named here only so the gap is explained;
%     nothing is inferred from them.
%   - Pages 5, 7, 9 and 15 are four leaves (« -0.13- », « -0.14- »,
%     « -0.15- », « -0.18- ») of a typed English text on quotients of a
%     scheme by a group action (categorical and geometric quotients, free
%     and proper actions, Definitions 0.6-0.8, Proposition 0.9, Example 0.4,
%     Lemma 0.5), not by him and unnamed on the leaves; 5, 7, 9 and 13 were
%     scanned upside down, being versos of his sheets. Each carries pencil
%     annotations — two of them substantial marginal notes on page 5 — and
%     is therefore kept, per the rule of folder 29 (hand.md §6): the typed
%     text is summarised in a French \note (not recomposed, the text not
%     being his), his pencil in \marginal.
%
% What the batch is. One manuscript run in blue ink, his own numbering
% « 1 », « 1 bis », « 2) », « 3 » on pages 2, 3, 4, 6, then unnumbered
% leaves 8, 10, 12, 14, 16, 20 and a leaf « 2 bis » on buff paper (18).
%   2      Theorem: S loc. noeth., U open, T its closed complement with all
%          residue characteristics = p; the functor F: A ↦ (A_U, T_p(A),
%          id) from abelian schemes over S to triples (A_U, M, φ), M a
%          p-divisible group; 1) U meets all components ⇒ F fully faithful;
%          2) under geometric unibranchness and a finiteness condition
%          (normalisation of S_red finite), F an equivalence; a struck
%          condition b) on strict localisations; a simplicial set
%          π_0(S') ⇇ π_0(S'') ⇚ π_0(S''') and a condition b').
%   3      Lemma (« 1 bis »): for u_ℓ: T_ℓ(A) → T_ℓ(B) there is a closed
%          subscheme T of S such that u_ℓ S' comes from u: A_S' → B_S' iff
%          S' → S factors through T (« Cf Invariants 1.2 »); Complément;
%          Corollary (a, b, c ⇒ u_ℓ algebraic) with its proof, cqfd.
%   4      « 2) »: the triples again; « Spécialisation »; 1) Foncteur
%          pleinement fidèle, via [ ] 1.2 and the schematic closure X_0;
%          a long struck Serre-Tate / completion argument; a crowded boxed
%          marginal block on counterexamples (two complete dvr spectra glued).
%   6      « 3 »: 2) Foncteur ess. surjectif; a) reduction to S reduced via
%          Serre-Tate, formal abelian schemes over T_n, S' = Spec Λ^, and
%          fpqc descent along U ⊔ S' → S (« d'après la réduction de
%          Raynaud »); b) struck.
%   8, 10  b) case S geom. unibranch via a Jacobi rigidification and a
%          modular scheme R, properness of S' → S, Main Theorem, Tate's
%          theorem; c) case S normal: level-ℓ rigidification (ℓ ≠ p, 3),
%          normalisation S' of S in R(U')/R(U), A' over S' with
%          T_p(A') ≃ p*(M), descent via S'' = S' ×_S S'.
%   12, 14 descent data on A' relative to S' → S (SGA 1 VIII), étale
%          splitting of _ℓ(A_U); d) case S reduced; a struck effectivity
%          attempt; e) general case, gluing the local solutions A(s) on
%          S(s) = Spec O_{S,s}.
%   16, 18 A ×_S U ≃ A_U and T_p(A) ≃ M; a Lemma on « systèmes locaux » of
%          abelian schemes (with his « Douteux » in the margin); page 18's
%          Lemma on extending an isomorphism given at Spec O_{S,t} to a
%          neighbourhood.
%   20     The Question « quand u_p provient-il d'un hom u: A → B », the
%          largest closed subscheme T on which u_p is algebraic, and a
%          Corollary via the set Z of associated and maximal points — a
%          restatement of page 3's lemma with ℓ replaced by p.
% The ordering of 12 → 14 → 16 and the attachment of 18 (« 2 bis ») follow
% the continuations the sentences show where they show one; where they do
% not, only the page order is given.
%
% Hand. Fast and heavily reworked, with whole paragraphs framed and struck
% by diagonals (pp. 2, 3, 4, 6, 8, 14, 20). Statements and formulas carry;
% much connective prose does not and is left \ill{}. Large struck blocks
% are given as one \struck{} with one \note, per hand.md §6.
%
% Notation fixed for the batch. T_p(A) for his Tate module / p-divisible
% group of A (he writes T_p for both); on page 4 only, T_p carries a small
% dark superscript mark, transcribed \natural and flagged (it may be a
% small φ or a completion mark); T_ℓ, u_ℓ, u_p as written; {}_ℓ(A_U) for his
% left-subscript ℓ (the ℓ-torsion); A_U, A_{S'}, A_{X_0} for base changes;
% S_red, Spec Λ, Λ^ (\hat\Lambda), T_n, T_0 as he writes them; Ass O_S;
% « car. k(t) » kept; « loc. noeth. », « pl. fid. », « ess. surjectif »,
% « géom. unib. », « s.f. », « chgt de base », « tte comp. conn. » kept.
% His empty square brackets « [ ] 1.1 », « [ ] 1.2 », « [ ] p. 75-76 » are
% his, a reference left to be filled, and are kept as « [\;] ». Diagrams
% sketched in the margins are set as arrays where their arrows are not
% certain enough for tikz-cd.
%
% Cross-references out of the batch: none explicit to pages 21-52. The
% argument of page 16 breaks off mid-sentence; page 18 (« 2 bis ») stops
% on « Alors » and page 20 is self-contained.
\documentclass[11pt,a4paper]{article}
% Le préambule commun à toutes les transcriptions du fonds Grothendieck.
%
% Il tient en une idée : **l'incertitude se compose, elle ne se résout pas.**
% Une transcription qui lisse un mot douteux a détruit la seule chose pour
% laquelle on la faisait — un lecteur doit pouvoir savoir, à chaque endroit,
% ce qui était sur le feuillet et ce qui a été ajouté. Les six macros qui
% suivent sont donc l'appareil critique, et non de la décoration.
%
% Les mêmes commandes sont reconnues par `scripts/render.mjs`, qui produit la
% vue de lecture du site. Une macro ajoutée ici doit l'être aussi là-bas,
% faute de quoi le rendu échouera bruyamment — ce qui est voulu : un rendu
% silencieusement faux est pire qu'un rendu absent.

\ProvidesPackage{grothendieck}[2026/08/08 Transcriptions du fonds Grothendieck]

\usepackage{amsmath,amssymb,amsthm}
\usepackage{mathrsfs}
\usepackage{stmaryrd}
% Les diagrammes commutatifs sont partout dans ce fonds : c'est la langue
% même de la géométrie algébrique de Grothendieck, et une transcription qui
% les rendrait en prose ne transcrirait rien.
\usepackage{tikz-cd}
% Français par défaut — c'est la langue des feuillets — mais l'anglais est
% chargé aussi : la traduction et le résumé sont des documents anglais, et la
% typographie française (espaces avant les deux-points) y serait une faute.
% Ces documents basculent par \selectlanguage{english} en tête de corps.
\usepackage[english,french]{babel}
\usepackage{fontspec}
\usepackage{geometry}
\geometry{a4paper,margin=25mm,marginparwidth=28mm}
\usepackage{marginnote}
\usepackage{xcolor}
% ulem, not soul. `soul`'s \st and \ul are built on a character-by-character
% scan that XeLaTeX's font machinery breaks: the document fails with an
% undefined control sequence rather than degrading. ulem does the same two jobs
% and is Unicode-engine safe. `normalem` stops it hijacking \emph, which the
% transcriptions use for Grothendieck's own emphasis.
\usepackage[normalem]{ulem}
\usepackage{hyperref}
\hypersetup{colorlinks=true,linkcolor=black,urlcolor=black,pdfborder={0 0 0}}

% --- Métadonnées du lot -----------------------------------------------------
% Reprises telles quelles par le script de rendu, qui les affiche en tête de
% la vue de lecture. Elles fixent ce que le fichier prétend couvrir : sans
% elles, une transcription flotte sans point d'attache dans un fonds de seize
% mille feuillets.

\newcommand{\folder}[1]{\gdef\grothFolder{#1}}
\newcommand{\batch}[1]{\gdef\grothBatch{#1}}
\newcommand{\leaves}[2]{\gdef\grothFirst{#1}\gdef\grothLast{#2}}
\newcommand{\foldertitle}[1]{\gdef\grothFoldertitle{#1}}
\gdef\grothFolder{?}\gdef\grothBatch{?}\gdef\grothFirst{?}\gdef\grothLast{?}\gdef\grothFoldertitle{}

% --- Le résumé --------------------------------------------------------------
%
% Il ouvre la lecture modernisée, sous le titre « Résumé », et il est composé
% en petit. Ce n'est pas un ornement : c'est le seul passage du document qui
% ne soit pas des mathématiques, et un lecteur doit pouvoir le reconnaître
% avant de l'avoir lu — puis le sauter s'il connaît déjà le sujet, ou s'y
% arrêter s'il ne le connaît pas. Le filet qui le suit marque la reprise du
% corps.

\newenvironment{resume}
  {\small\setlength{\parskip}{0.45em}}
  {\par\medskip\hrule\medskip}

% --- Mots-clés ---------------------------------------------------------------
%
% En anglais, en clôture du résumé de la lecture modernisée : le vocabulaire
% moderne sous lequel chercher ce que les feuillets construisent. Le manifeste
% du site (`npm run manifest`) les extrait pour étiqueter la cote entière —
% cette ligne est la source unique des tags, il n'y a pas de fichier à côté.
\newcommand{\keywords}[1]{\par\smallskip\noindent
  {\footnotesize\textsc{Keywords} --- \textit{#1}}\par}

% --- L'appareil critique ----------------------------------------------------

% Le numéro de feuillet, en marge. C'est l'ancre qui permet de régler un
% désaccord sur une formule en regardant *un* feuillet plutôt qu'une chemise
% de six cents. Le site s'en sert aussi pour tourner les pages du fac-similé
% quand on fait défiler la transcription.
\newcommand{\leaf}[1]{%
  \marginnote{\footnotesize\color{gray!70}\textsf{#1}}%
  \label{leaf:#1}\ignorespaces}

% Illisible : on ne devine pas. Un mot inventé qui se lit comme les autres est
% le pire résultat possible.
\newcommand{\ill}{\textcolor{red!60!black}{[\,\ldots\,]}}

% Lecture douteuse : le mot est donné, et signalé comme incertain.
\newcommand{\uncertain}[1]{\uline{#1}}

% Ajout de l'éditeur — une lettre manquante, un mot sous-entendu.
\newcommand{\add}[1]{\textcolor{blue!50!black}{[#1]}}

% Biffé par Grothendieck lui-même. Ce qu'il a rayé fait partie du document :
% une rature dit souvent où la pensée a changé de direction.
\newcommand{\struck}[1]{\sout{#1}}

% Note du transcripteur, sur l'état matériel du feuillet.
\newcommand{\note}[1]{\par\smallskip{\small\color{gray!60!black}\textit{[#1]}}\par\smallskip}

% Note marginale de Grothendieck, distincte du corps du texte.
\newcommand{\marginal}[1]{\par\smallskip\noindent\hspace{1em}%
  {\small\color{gray!60!black}#1}\par\smallskip}

% --- Titre ------------------------------------------------------------------

\AtBeginDocument{%
  \begin{center}
    {\large\bfseries Fonds Alexandre Grothendieck --- cote n\textsuperscript{o}~\grothFolder}\\[2pt]
    {\itshape\grothFoldertitle}\\[4pt]
    {\small feuillets \grothFirst--\grothLast\ (lot \grothBatch)}\\[2pt]
    {\footnotesize\color{gray!60!black}Transcription établie d'après le fac-similé numérisé
     par l'Université de Montpellier.\\
     Les lectures incertaines sont soulignées, les illisibles notées [\,\ldots\,],
     les ajouts entre crochets.}
  \end{center}
  \vspace{1em}}

\endinput


\folder{100}
\batch{1}
\pages{1}{20}
\dating{[à partir de 1967-1968]}
\watermark{Édition de démonstration}
\foldertitle{Bonne réduction des variétés abéliennes via bonne réduction de Barsotti-Tate Tp (A) (base quelconque) : notes manuscrites (s.d.), copies de tapuscrit annoté (s.d.).}

\begin{document}

\section{Bonne réduction et groupes $p$-divisibles}
\note{titre de sa main, en haut à droite du feuillet 1, une chemise de
papier bistre par ailleurs vierge qui ne reçoit pas de numéro de page}

\page{2}
\note{Grothendieck numérote ce feuillet « 1 » ; le suivant « 1 bis », puis
« 2 », « 3 »… Écriture rapide, surchargée de reprises ; une large partie
centrale est encadrée et biffée de traits obliques.}

\textbf{Théorème.} $S$ schéma loc.\ noeth., $U$ ouvert, $T$ fermé
complémentaire, $p$ un premier, tel que $\forall t \in T$, car $k(t) = p$.
Considérons la catégorie $C$ des schémas abéliens $A$ sur $S$, et la
catégorie $C'$ des couples $(A_U, M, \varphi)$, où $A_U$ est un schéma
abélien sur $U$, $M$ un groupe $p$-divisible sur $S$, et
$\varphi : M|U \to T_p(A_U)$. Soit
\[
  F : C \longrightarrow C'
\]
le foncteur
\[
  F(A) = \bigl(A \times_S U,\ T_p(A),\ \mathrm{id}_{T_p(A)|U}\bigr).
\]
\note{le mot « couples » désigne ici des triples ; l'exposant « $A$ » est
ajouté au-dessus de « schémas abéliens », pour nommer l'objet}

Alors :

1) Supposons que $U$ \uncertain{rencontre toutes les composantes
\uncertain{connexes}} de $S$ (i.e.\ $U$ dense dans $S$, i.e.\
$\mathrm{prof}_T(S) \geqslant 1$). Alors le foncteur $F$ est pleinement
fidèle.
\note{la parenthèse « (i.e.\ $U$ dense dans $S$, i.e.\
$\mathrm{prof}_T(S)\geqslant 1$) » et la ligne « $U$ rencontre toutes
\ldots » sont ajoutées d'une encre plus claire au-dessus de la ligne ; le
dernier mot est mal formé}

\struck{\ill{}) $\exists$ un morphisme fini surjectif $S' \to S$ \ill{},
et que $U$ soit dense dans $S'$ et que \ill{}}

2) Supposons \struck{que} \ill{} (i.e.\ le normalisé de $S_{\mathrm{red}}$
fini sur $S$, \uncertain{si} $S$ local).
\note{le mot qui suit « Supposons » est recouvert ; au début de la ligne
suivante, « $S'$ géom.\ unibranche », peut-être la condition annoncée}

\struck{b) $\forall\, t \in T$, si $S'$ est le localisé strict de $S$ en
$t$, $\bar t$ son point fermé, et $S'_1$ et $S'_2$ deux composantes
irréductibles de $S'$, on ait $S'_1 \cap S'_2 \neq \{\bar t\}$. (N.B.\ cette
condition \uncertain{signifie aussi que} \ill{} $F$ est \ill{} image inverse
$U''$ de $U$ dans $S'' = S' \times_S S'$ \ill{} \uncertain{catégories}
\ill{} $S''$. Elle est \ill{} des points de $T$).}
\note{tout ce paragraphe b) est encadré et biffé de trois traits obliques ;
en marge, dans le cadre, une reprise « b) Posant », elle aussi enfermée, et
« b) Comme. », que rien ne suit}

Alors $F$ est une équivalence de catégories \struck{pleinement fidèle}.
\note{sous le cadre, à gauche, isolé : « $\mathrm{prof}_T(S) \geqslant 2$ »,
puis le mot « catégories »}

\marginal{Pour 1) TSVP}
\note{cerclé, avec une flèche vers le bord droit}

Soit, avec les notations de a), $S'' = S' \times_S S'$,
$S''' = S' \times_S S' \times_S S'$. Considérons l'ensemble simplicial
\[
  \underset{L_0}{\pi_0(S')} \leftleftarrows \underset{L_1}{\pi_0(S'')}
  \;\substack{\leftarrow\\[-2pt]\leftarrow\\[-2pt]\leftarrow}\;
  \underset{L_2}{\pi_0(S''')},
\]
\struck{Com} et soit
\[
  K_0 \leftleftarrows K_1 \;\substack{\leftarrow\\[-2pt]\leftarrow\\[-2pt]\leftarrow}\; K_2
\]
le sous-ensemble simplicial formé avec les composantes connexes dont l'image
dans $S$ rencontre $U$. Supposons que : \struck{pour tout}
b$'$) la plus petite partie de $L_1$ contenant $K_1$ et st.\ stable par
« \uncertain{troisième face} » soit $L_1$ (\ill{} par la condition b)).
Alors $F$ est \struck{pl.\ f.} une équivalence de catégories\,\ldots
\note{les noms $L_0, L_1, L_2$ sont écrits au-dessus des $\pi_0$ ; au bas du
feuillet, sous la dernière ligne, une incise d'une lecture incertaine,
« ce qui est \ill{} », renvoie à la condition b)}

\page{3}
\note{numéroté par lui « 1 bis »}

\textbf{Lemme.} $A$, $B$ schémas abéliens sur $S$ loc.\ noeth.,
$u_\ell : T_\ell(A) \to T_\ell(B)$. Alors $\exists$ sous-schéma fermé $T$ de
$S$ tel que, pour tout $S'/S$, on ait : $u_{\ell S'}$ provient de
$u : A_{S'} \to B_{S'}$ ssi $S' \to S$ se factorise par $T$.
\note{« on ait : $u_{\ell S'}$ » est ajouté en marge gauche, avec un trait
de renvoi}

Cf Invariants 1.2 (\ill{} démonstration \uncertain{analogue}).

\textbf{Cor} \struck{Supposons que $\forall s \in \mathrm{Ass}$ \ill{}
l'image de \ill{}}

\textbf{Complément.} De plus, \struck{\ill{}} pour tout $t \in T$ tel que
car.\ $k(t) = \ell$, $T$ \uncertain{majore} un voisinage ouvert de $t$.
\note{au-dessus de la rature, en interligne, deux mots illisibles}

\textbf{Corollaire.} Supposons que l'ouvert $U \subset S$ soit tel que
\begin{itemize}
  \item[a)] tte comp.\ conn.\ de $S$ rencontre $U$,
  \item[b)] $t \in S - U \Rightarrow$ car.\ $k(t) = \ell$,
  \item[c)] $u_\ell|U$ provient de $u_U : A|U \to B|U$.
\end{itemize}
Alors $u_\ell$ provient de $u : A \to B$.

\struck{En effet, \ill{} on peut supposer $S$ connexe, \ill{} irréductibles
\ill{} \ill{} distinguons $t \in T \cap (S-U)$ \ill{} Comme $T \neq
\emptyset$, \ill{} on aurait une comp.\ \ill{} de $S$ telle \ill{} que \ill{}
comp.\ irréd.\ de $S$ qui rencontre $T$, $\bar U$ \ill{} $T$ \ill{} $\in T$.
Or \ill{} $S_i \cap U = \emptyset$ \ill{} \uncertain{complément}. Donc $T
\supset S_{\mathrm{red}}$.}
\note{tout ce passage est encadré et biffé de deux longs traits obliques ;
en marge du cadre, une incise cerclée « \uncertain{grâce à} b) », et
« $\Rightarrow U \neq \emptyset$ »}

Mais $\underline{T}$ est \struck{fermé} ouvert : $T$ \uncertain{est ouvert},
\uncertain{contient} $U$ \struck{et} par hyp.\ c), et sur les pts de
$(S - U) \cap T$ grâce au complément et de b). Donc $\overline{T}$ est un
sous-schéma ouvert, et comme il est fermé, c'est une réunion de comp.\
connexes. Par a), on a donc $\overline{T} = S$,

cqfd.
\note{le $T$ de « Mais $T$ » est souligné, celui de « Donc $T$ » et de
« $T = S$ » surmonté d'un trait ; « sous-schéma » est ajouté au-dessus
d'« ouvert »}

\page{4}
\note{numéroté par lui « 2) » ; en tête, un signe isolé en forme de croix
surmontée, peut-être un renvoi}

\[
  T \subset S \supset U = S - T
\]
$T$ : car.\ résiduelles toutes $= p > 0$.
\[
  A \longmapsto \bigl(A_U,\ T_p^{\natural}(A),\ \varphi\bigr)
\]
schémas abéliens sur $S$ $\longmapsto$ triples $(A_U, M, \varphi)$, $A_U$
schéma abélien sur $U$, $M$ \uncertain{gpe $p$-div.}\ sur $T$,
$\varphi : T_p^{\natural}(A_U) \simeq M$.
\note{l'exposant de $T_p$, ici noté $\natural$, est un petit signe ajouté
d'une encre plus foncée, un rond plein traversé d'un trait (il pourrait
s'agir d'un $\varphi$ minuscule ou d'une marque de complétion) ; il revient
à chaque occurrence de $T_p$ sur ce feuillet}

\textbf{Spécialisation}, $U$ \struck{dense} \struck{schématiquement}
\struck{schématiquement} dense et $S - U$ de car.\ $p$.
\marginal{Pour pleine fidélité, \struck{\ill{}} \uncertain{suffit} de
$U$ dense \ill{} \struck{\ill{}} (\ill{} rencontre $\bar U$ \ill{} p.\ ex.\
$S$ \ill{}) \struck{\ill{} fidélité \ill{}} \ill{} contre-exemple \ill{}.
On \uncertain{obtient} \uncertain{en recollant} deux spectres d'anneaux de
v.d.\ discrète complets, et $U =$ un des deux pts
\uncertain{maximaux}\,\ldots\,) Y a-t-il des contre-exemples avec $S$
\ill{} dim.\ $2$ ?}
\note{ce bloc, à droite en haut du feuillet, est cerné d'un trait et
lourdement repris : cinq ou six lignes y sont biffées à l'encre épaisse et
ne se lisent pas}

1) \textbf{Foncteur pleinement fidèle}

Fidèle bien connu, car $A \mapsto T_p^{\natural}(A)$ est fidèle.
« Pl. » fidèle : il faut prouver que si
$u_p : T_p^{\natural}(A) \to T_p^{\natural}(B)$ \uncertain{est algébrique}
sur $U$, il l'est sur $S$. Si $U$ est sch.\ dense dans $S$, cela résulte de
[\;] 1.2. Dans le cas général, soit $X_0$ l'adhérence schématique de $U$ dans
$X$. Alors on dispose d'un homomorphisme
\[
  \varphi_0 : A_{X_0} \to B_{X_0},
\]
\uncertain{compatible avec} $u_p$. \struck{Comme}
\note{« compatible avec $u_p$ » est ajouté en interligne ; les crochets
vides « [\;] » sont les siens, une référence à compléter}

\struck{Soit $\mathcal{J}$ l'idéal qui définit $X_0$ sur $X$, il est
\uncertain{nilpotent} \ill{} $T$. Alors le th.\ de Serre-Tate implique que
\ill{} la restriction de $u_0$ à \ill{} ($= T$ \ill{} réduit) \ill{}
$A_{T_n} \to B_{T_n}$, où $T_n$ est le $n$-ième voisinage infinitésimal
de $T_0$ dans $S$. Se ramenant au cas où $S = \mathrm{Spec}\,\Lambda$
\ill{}, \ill{} un homom.\ $A_{\hat\Lambda} \to B_{\hat\Lambda}$,
$\hat\Lambda = $ complété de $\Lambda$ pour l'idéal de $T_0$ dans $S$.
Comme \ill{} fid.\ plat, \ill{} $U \times \mathrm{Spec}\,\hat\Lambda \to
\mathrm{Spec}\,\Lambda$ \ill{}}
\note{ce long paragraphe est encadré et biffé de deux traits obliques ; hors
du cadre, en bas à droite, reste non biffé : « car \uncertain{il est} \ill{}
[\;] 1.1. »}

\page{5}
\note{Feuillet de tapuscrit anglais, paginé « -0.13- », dont le feuillet 4
occupe le verso ; ce tapuscrit n'est pas de lui et le feuillet ne le nomme
pas. Il n'est pas recomposé ici : le texte dactylographié est résumé en
français, et seules ses annotations au crayon sont transcrites. Le feuillet
contient la fin d'une démonstration (surjectivité de $\phi'$, puis
« by Zariski's Main Theorem, $\bar\omega$ is an isomorphism »), une
« Remark (5) » donnant trois conditions équivalentes pour que $(Y,\phi)$
soit un quotient catégorique universel de $X$ par $G$, et un alinéa « (6) »
énonçant trois conditions suffisantes pour que $(Y,\phi)$ soit un quotient
catégorique : i) $\phi\circ\sigma = \phi\circ p_2$ ; ii)
$\mathcal{O}_Y$ est le sous-faisceau des invariants de
$\phi_*(\mathcal{O}_X)$ ; iii) l'image par $\phi$ d'un fermé invariant est
fermée, et $\phi(\bigcap W_i) = \bigcap \phi(W_i)$ pour toute famille de
fermés invariants.}

\note{en marge, un « ! » au crayon devant « by Zariski's Main Theorem »,
souligné d'un trait ondulé}

\marginal{$R \underset{p_2}{\overset{p_1}{\rightrightarrows}}$ \quad
$\varphi : X \xrightarrow{f} Y$ morphism of preschemes
\struck{\ill{}} $f p_1 = f p_2$ \uncertain{we say} $Y$ is a quotient $X/R$
if (i) $\mathcal{O}_Y$ is the sheaf of invariants of \ill{}
$\mathcal{O}_Y \to f_*(\mathcal{O}_X) \rightrightarrows
g_*(\mathcal{O}_R)$ \uncertain{is exact} (ii) Any invariant open $U$ of $X$
is the \uncertain{inverse image} of an \uncertain{open subset} of $Y$}
\note{note au crayon, en anglais, écrite en oblique dans la marge gauche,
en face de « (6) » ; l'indice de $f$ et $g$ est d'une lecture incertaine}

\marginal{I \uncertain{think} sufficient if $\varphi$ is qu.\ cpct qu.\
sep.}
\note{au crayon, au pied du feuillet, sous la condition iii), qu'une
accolade au crayon embrasse ; « qu.\ cpct » est souligné}

\page{6}
\note{numéroté par lui « 3 »}

2) \textbf{Foncteur ess.\ surjectif}
\marginal{« Si $S$ géom.\ unibranche en les points de $S - U$ (\ill{}) et
si le normalisé de $S_{\mathrm{red}}$ est fini sur $S$. »}
\note{encadré en haut à droite, relié par des guillemets doubles au titre ;
un trait le relie aussi à la parenthèse ci-dessous}

\struck{a) Supposons}

Partons de $A_U, M, \varphi$, il faut construire $A$. On peut supposer
\struck{$A$} \struck{\ill{}} $S = \mathrm{Spec}\,\Lambda$.
\struck{Par descente, il suffit de \ill{}}
[plus généralement, si, posant $S'' = S' \times_S S'$,
$U' = U \times_S S'$, $U'' = U \times_S S''$, $U'$ ($U''$) rencontre
toutes les composantes connexes de $S'$ ($S''$)]

a) Réduction au cas $S$ réduit.

Soit $S_0 = S_{\mathrm{red}}$, $A_0$ schéma abélien sur $S_0$ qui y vient
\ill{} de $M$. Par Serre-Tate, on trouve \struck{des} un schéma formel
abélien sur les $T_n$, donc un schéma abélien $A'$ sur
$S' = \mathrm{Spec}\,\hat\Lambda$. Son image inverse sur $S'_0$ est un vrai
schéma abélien, \struck{\ill{}} savoir l'image inverse de $A_0$ sur $S'_0$.
Donc (réf.\ \ldots\ !!!), $A'$ provient d'un vrai schéma abélien $A'$ sur
$S'$ [l'immersion $S'_0 \to S'$ \uncertain{étant} \uncertain{nilpotente}].
Or
\[
  U \sqcup S' \longrightarrow S
\]
est fid.\ plat, et on dispose d'une donnée de descente sur
$A_U \sqcup A'$. Elle est effective, \uncertain{puisqu'elle} dépend de
$S_0 \to S$, donc elle l'était déjà avant, d'après la réduction de
Raynaud. \uncertain{Ce qu'on gagne} !
\note{dans la marge gauche, en face de « $S' = \mathrm{Spec}\,\hat\Lambda$ »,
un petit carré :}

\begin{tikzcd}[column sep=small, row sep=small]
  S_0 \arrow[d] & S'_0 \arrow[l] \arrow[d] \\
  S & S' \arrow[l]
\end{tikzcd}

\struck{b) Réduction au cas $S$ local, \ill{} $T$ \uncertain{réduit} à un
point fermé.}

\struck{Pour faire cette réduction, on est ramené à prouver ceci : si
$t \in T$, $V$ un ouvert \ill{} contenant $t$ et $U$, \struck{et} $A_V$ un
schéma abélien sur $V$ \struck{prolongeant} \ill{} $M$ \uncertain{que} se
\ill{} $\mathcal{K}$, alors $A_V \to$ \ill{} ; $S' = \mathrm{Spec}\,
\mathcal{O}_{S,t}$ \ill{} voisinage, i.e.\ l'isom.\ donné
$T_p(A_{S'}) \simeq M_{S'}$ se prolonge en un isom.\ en un \ill{}}
\note{ce paragraphe b) est encadré et biffé de deux traits obliques ; un
« ? » au crayon dans la marge gauche}

\page{7}
\note{Même tapuscrit, feuillet « -0.14- », résumé : démonstration du
critère de l'alinéa (6) (on montre que $\phi$ est surjectif, puis, pour un
$S$-morphisme $\psi : X \to Z$ avec $\psi\circ\sigma = \psi\circ p_2$, on
construit un recouvrement ouvert $\{U_i\}$ de $Y$ tel que
$\psi^{-1}(V_i) \supset \phi^{-1}(U_i)$) ; puis l'alinéa « (7) » : les
conditions (i), (ii), (iii) de la Définition 0.6 sont conservées par tout
changement de base $Y' \to Y$, la condition (iv) seulement si $Y' \to Y$
est plat, par l'exactitude de la suite
$0 \to \mathcal{O}_Y \to \phi_*(\mathcal{O}_X)
\xrightarrow{\sigma^* - p_2^*} \psi_*(\mathcal{O}_{G\times_S X})$ et la
commutation de l'image directe au changement de base plat.}

\marginal{$\psi$}
\note{au crayon, en marge de « Now suppose $\phi : X \to Z$ », où le
$\phi$ dactylographié est barré au crayon : la correction en $\psi$ ; au
même crayon, « surjective » est souligné et
« $\psi^{-1}(V_i) \supset \phi^{-1}(U_i)$ » est entouré}

\marginal{$\varphi$ quasi-compact \uncertain{\&} quasi-separated \&}
\note{au crayon, en oblique dans la marge gauche, avec un trait de renvoi
devant « is flat » dans « However, if $Y' \to Y$ is flat, then condition
(iv) is preserved »}

\page{8}
\note{sans numéro de sa main ; la suite du feuillet 6}

\struck{voisinage ouvert $W$ de $t$, \ill{} ce qui \ill{} donné sur
$W \cap U$. En fait, \struck{d'après le th.\ de \ill{}} si on dispose
\ill{} de Raynaud \ill{} Tate \ill{} \ill{} \ill{} incidence canonique,
\ill{} \ill{} pouvant prendre}
\note{les cinq premières lignes du feuillet, qui terminent le paragraphe
b) biffé du feuillet 6, sont encadrées et biffées de traits obliques}

b) Cas $S$ \struck{\ill{}} \uncertain{géom.\ unib.}, $A_U$ \ill{}
d'une rigidification de Jacobi. On a donc un morphisme
\[
  U \longrightarrow R
\]
$R$ schéma \uncertain{modulaire} à décalage [utilisant aussi une
polarisation de $A_U$], d'où par le \uncertain{critère} \uncertain{valuatif}
\uncertain{de l'adhérence des} \ill{} \uncertain{propre} un
\uncertain{sous-schéma} fermé $S'$ de $S \times R$. Le critère
\uncertain{valuatif} et le résultat de Raynaud dans le \ill{} \uncertain{ou}
d'un \uncertain{trait} \uncertain{montrent} \uncertain{que} $S' \to S$ est
propre. Je veux : prouver que c'est un isom.\ [je reviens au th., grâce à
Tate, à prouver que $A_U$ se prolonge en $A$, quand $T_p(A_U)$ se prolonge
en un $M$ et $M$ est \ill{} ; \ill{} unique à iso.\ unique
près]. Par le Main Theorem, \ill{} on va \uncertain{prouver} que les fibres
$S'_s$ sont discrètes. \struck{les pts de $S'$} Introduisons
\ill{} $S''$ de $S'$, \uncertain{qui} \uncertain{normalisé} d'abord
\uncertain{fini} sur $S'$. Alors le th.\ de Tate implique que la
\uncertain{bonne réduction} $T_p(A_{S''})$ \uncertain{image inverse de}
$A_R$ \ill{} \struck{$T_p(A)$} $M$. D\ill{} la restriction de
$A_{S''}|S''_s$, \uncertain{correspondant} \ill{} $T_p$ « constant », et
$S''_s$ étant géom.\ connexe,
\note{le raisonnement se poursuit au feuillet 10}

\[
  \begin{array}{ccc}
    & S'' & \hookleftarrow\; U'' \\
    & \downarrow & \wr \\
    R \longleftarrow & S' & \hookleftarrow\; U' \\
    & \downarrow & \wr \\
    & S & \hookleftarrow\; U
  \end{array}
\]

\note{croquis de la marge gauche : $S'' \leftarrow U''$ au-dessus de
$S' \leftarrow U'$ au-dessus de $S \leftarrow U$, les flèches verticales
à gauche tracées vers le bas, les verticales de droite marquées « $\wr$ »,
et une flèche de $S'$ vers $R$ ; le sens des flèches verticales de gauche
est incertain}

\page{9}
\note{Même tapuscrit, feuillet « -0.15- », résumé : fin de l'alinéa (7)
(changement de base plat, et quotients géométriques universels), une
« Definition 0.7 » : une action $\sigma$ de $G$ sur $X$ est libre si
$\Psi = (\sigma, p_2) : G \times_S X \to X \times_S X$ est une immersion
fermée ; puis la « Definition 0.8 » de la variété de Picard et la
« Proposition 0.9 » (Cartier, Nishi) : un $S$-préschéma en groupes
$G$ de type fini sur un corps est lisse en caractéristique $0$.}

\marginal{Dire : \emph{propre} \uncertain{et libre}}
\note{au crayon, en oblique dans la marge gauche en face de la
« Definition 0.7 » ; au-dessus, d'une écriture plus pâle, trois mots,
lus \uncertain{compliqué des termes vagues} ; en bas à droite, le
« $S$ » de « $e_X : X \to S$ » est repassé à l'encre}

\page{10}
\note{sans numéro de sa main ; les six premières lignes sont barrées d'un
trait oblique}

\struck{\ill{} coïncident (cf lemme \ill{}) que $A_{S''}$ est \uncertain{constant}.
Il \uncertain{reste} \ill{} \ill{} que \uncertain{pourvu} \uncertain{que}
\uncertain{polarisé}, \ill{} \uncertain{finitude}. Il \ill{} \ill{} le choix
que d'un \ill{} fini de rigidification de Jacobi, \ill{} $(S''_s)_{\mathrm{red}}
\to R$ \ill{} factorise par \uncertain{$k(s)$}, donc $S''_s$ est
\uncertain{radiciel} \ill{} \uncertain{bien}. Ce qu'on gagne.}

Le cas où \ill{} un peu plus \uncertain{général} $S''$ fini sur $S'$ se
traite par la méthode essentiellement de [\;] p.\ 75--76,
\uncertain{condition} \uncertain{qui} \ill{} \ill{} \ill{} que l'image
\uncertain{réciproque} sur $U'' \simeq U\,$\ill{} \ill{} $T_p(A'')|U''$
et $M_{S''}|U''$ se prolonge \ill{} $S''$. \struck{Il faudrait}
\note{au-dessus de « $T_p(A'')|U''$ et », en interligne : « o.k » ;
« essentiellement de [\;] p.\ 75--76 » est souligné}

\struck{\ill{} \ill{} certains \ill{} quantités,}

c) Cas $S$ \struck{\ill{}} \uncertain{normal} \uncertain{unibranche}.
Choisissons un \ill{} fini \ill{} $U' \to U$ un \uncertain{morphisme}
\ill{} $\ell \neq p, 3$, et \uncertain{sur} \ill{} étale \uncertain{et}
surjectif \struck{\ill{}} tel que $A_{U'}$ soit
\uncertain{muni} d'une rigidification de Jacobi d'échelon $\ell$.
\ill{} Normalisons $S$ dans $R(U')/R(U)$, d'où \uncertain{$R''(S)$}
un \uncertain{morphisme} fini (\,$S' \to S$), qui \uncertain{prolonge}
$U' \to U$, \ill{} $S'$ normal intègre.
\marginal{\uncertain{inutile} \uncertain{au} \uncertain{fond} !}
\note{cerclé dans la marge gauche et relié par une boucle à
« Normalisons » ; « $R''(S)$ » et « prolongeant $U' \to U$ » sont
entourés en interligne}

Alors on peut trouver un schéma abélien $A'$ sur $S'$ qui prolonge
$(A_U) \times_U U'$, avec $T_p(A') \simeq p^{*}(M)$. Alors par le
\uncertain{critère} « pl.\ fid. » appliqué à
$S'' = S' \times_S S'$, utilisant que $U''$ est dense dans $S''$
(\struck{\ill{}} N.B.\ $S'' \to S$ est ouvert car $S' \to S$ \uncertain{est
universellement} ouvert\,\ldots), on trouve
\note{la phrase se poursuit au feuillet suivant}

\[
  \begin{array}{ccccc}
    U & \longleftarrow & U' & \leftleftarrows & U'' \\
    \cap & & \cap & & \cap \\
    S & \xleftarrow{\ p\ } & S' & \leftleftarrows & S''
  \end{array}
\]

\note{croquis de la marge gauche, redessiné ; les doubles flèches sont
les deux projections}

\page{12}
\note{sans numéro de sa main ; suite du feuillet 10}

que $A'$ est muni d'une \emph{donnée de descente} relativement à
$S' \to S$. En particulier, \struck{E} \uncertain{si} $A'$ est muni d'une
telle donnée de descente, \struck{que} \uncertain{comme} $S' \to S$ est
\uncertain{un morphisme} de descente effective pour les schémas abéliens
\uncertain{projectifs} finis étales (SGA 1 VIII), donc on trouve un
\struck{tel} tel revêtement \struck{R} \uncertain{en groupes} $N$ de $S$,
tel que
\[
  N|U \simeq {}_{\ell}(A_U).
\]
On montre que l'on peut trouver un morphisme \uncertain{fini}
\emph{étale} surjectif $S_1 \to S$ qui splitte ${}_{\ell}(A_U)$,
\struck{Remplaçant alors $S'$} et on peut supposer $S'$ \struck{fini}
\emph{étale} sur $S$, donc plat. Alors ($A'$ étant projectif sur $S'$) on
\struck{\ill{}} \struck{\ill{}} peut \emph{descendre} $A'$ en $A$, qui
répond à la \uncertain{q}.
\marginal{\struck{\ill{}} pourvu qu'il existe $S' \to S$ fini surjectif,
avec $S'$ géom.\ unib.}
\note{ajout encadré en interligne, relié à « répond » ; « descendre » est
souligné}

d) Cas $S$ réduit, ($S'$ fini sur $S$). Utilisant le fait que $S$ est
géom.\ unibranche en les points de $S - U$, on voit que l'image inverse
$U'' = U' \times_U U'$ de $U$ dans $S'' = S' \times_S S'$ est
\uncertain{dense} dans $S''$ : \struck{Par le cas b),} Par le cas c), on
trouve sur $S'$ un schéma abélien $A'$, et par la pleine fidélité (pour
$S''$, $U''$), on voit que $A'$ est muni d'une donnée de descente
relativement à $S' \to S$, \ill{} \ill{} \ill{} $S' \to S$. Notons que
un morphisme fini, \uncertain{et} un morphisme étant un épimorphisme fini,
\struck{E} est aussi de descente \uncertain{pour} les schémas abéliens, et
aussi pour les gpes $p$-divisibles. Donc on est ramené à
\note{la phrase se poursuit au feuillet 14}

\[
  \begin{array}{ccccc}
    U & \longleftarrow & U' & \leftleftarrows & U'' \\
    \cap & & \cap & & \cap \\
    S & \longleftarrow & S' & \leftleftarrows & S''
  \end{array}
\]

\note{croquis de la marge gauche, en face de « d) » ; dans le coin
inférieur gauche, deux mots de sa main lus
\uncertain{Remarques} \uncertain{digression}, séparés du texte par un
trait vertical}

\page{14}
\note{sans numéro de sa main ; suite du feuillet 12}

montrer que la donnée de descente envisagée est effective. Or on
\struck{est ramené facilement au cas où $S$ est local, de point fermé $s$,
$U = S - \{s\}$.} \struck{Par factorisation (où \uncertain{utilisant}}
\struck{l'hypothèse} \struck{en fait que $S'$ \uncertain{normal}) on est
ramené à un morphisme $S' \to S$ qui est un épimorphisme effectif}
\struck{(« $\mathrm{Spec}\,k' \to \mathrm{Spec}\,k$ »)} \struck{$R'/R$ de
deg.\ $1$ sur $R$ (?),} \struck{utilisant l'effectivité des descentes sur
les \ill{} de FGA} \struck{\ill{} on prouvera : montrer qu'il \ill{}}
\struck{est plat} \struck{\ill{} \uncertain{polarisation} \ill{} pour
\ill{} existe sur $A'$ \ill{}} \struck{\uncertain{donnée} de
\uncertain{descente}} \struck{une \uncertain{th.}\ de descente qui nous
\uncertain{suffira}, cf plus bas.}
\note{tout ce bloc, sept lignes, est encadré et biffé de huit traits
obliques ; à l'intérieur, « l'hypothèse », « $R'/R$ de deg.\ $1$ sur $R$
(?) », « utilisant l'effectivité … FGA », « est plat » et « donnée de
descente » sont en outre biffés chacun d'un trait propre ; dans la marge gauche, un îlot cerné où se lit
\uncertain{effective}, lui-même biffé ; sous le cadre, dans la marge, un
« ?! » barré d'un trait}

e) Cas général. \struck{Le cas} Cas \ill{} déjà \uncertain{vu} \ill{}
\uncertain{réduit} du d). \uncertain{Supposons} $S$ réduit. Le cas
$S$ \emph{local} \ill{}, \uncertain{on} trouve pour tout $s \in S$, de
trouver un schéma abélien $A(s)$ sur $\mathrm{Spec}\,\mathcal{O}_{S,s}
= S(s)$, qui \uncertain{soit} \uncertain{une} solution du pb induit sur
$S(s)$ (avec $U(s) = U \times_S S(s)$, $T(s) = T \cap S(s)$,
\uncertain{$M(s)$}, \ldots). Par unicité, les $A(s)$ forment un
\uncertain{système} \uncertain{cohérent} de schémas abéliens, i.e.\ pour
$t$ générisation de $s$, on a
\[
  A(t) \simeq A(s) \times_{S(s)} S(t),
\]
avec \ill{} \ill{} \ill{} \uncertain{transitivité}. \struck{Supposons que}
On \ill{} \ill{} deux catégories des schémas \ill{} un foncteur
\uncertain{pl.\ fid.} \uncertain{Mais} \ill{} la cat.\ des
\uncertain{systèmes} \uncertain{locaux} de \ill{} s.f.\ sur $S$ dans la
cat.\ des \ill{} s.f.\ sur les $S(s)$. Supposons qu'on \ill{} \ill{} de
s.f.\ sur les \ill{} \uncertain{naturelle}. \ill{} que $\{A(s)\}_s$
\uncertain{provient} d'un schéma \struck{abél.}\ $A$ sur $S$, \ill{}.
Alors on a \ill{} un schéma abélien \ldots\ (par la pleine fidélité
précédente).
\note{la marge gauche porte, sur toute la hauteur de ce paragraphe, une
longue note écrite en oblique, dont se lisent seulement « Lemme »
(souligné), « \uncertain{l'hyp.\ faite} », « $(A_U, M, \varphi)$ », « $T$ »
et, à la fin, « \uncertain{schémas abéliens} » ; le reste \ill{}}

\page{15}
\note{Même tapuscrit, feuillet « -0.18- », résumé : un « Example 0.4 »
d'un groupe semi-simple $G$ sur $\mathbb{C}$ agissant sur un schéma
quasi-affine non singulier de type fini, dont i) tous les stabilisateurs
sont triviaux (action ensemblistement libre), ii) un quotient géométrique
existe, qui est un schéma de type fini, d'où l'action est séparée, iii)
l'action n'est pas propre et en particulier pas libre au sens
algébro-géométrique ; début de la preuve avec $G = SL(2)$ agissant sur
$V_1 \times V_4$. Au crayon : « free » est entouré dans i), « not free »
est souligné d'un trait ondulé dans iii), et un trait marque la marge en
face de iii).}

\page{16}
\note{sans numéro de sa main ; ce feuillet ne suit pas directement le
feuillet 14}

$A \times_S U \simeq A_U$ car (\ill{} \ill{}). De plus, je dis qu'on a un
isom.\ $T_p(A) \simeq M$, qui induit les $T_p(A(s)) \simeq M(s)$. En
effet, par \ill{} \ill{} \ill{}, par le résultat déjà signalé de
pl.\ fid., on aura un isom.\ \struck{$T_p(A$} $p^{*}A \simeq M(s)$ qui
induit l'\ill{} les
\[
  \underset{\substack{\shortparallel\\ (p^{*}A)_s}}{p^{*}(A_s)}
  \;\simeq\;
  \underset{\substack{\shortparallel\\ (M(s))_s}}{M(s)(s)} .
\]
Sur $U$, c'est \ill{} \uncertain{bien} \uncertain{connu}
\uncertain{tautologique}. Donc \uncertain{tautologique}.

Il reste à prouver la pleine fidélité précédente
(\uncertain{\ill{}}), et le \ill{} \ill{}.

\textbf{Lemme} : Un « sys.\ local » de schémas abéliens sur \ill{}
provient d'un schéma [loc.\ noeth.] \ill{} un schéma abélien.
\marginal{Douteux, \ill{} \ill{}}
\marginal{[\;] \uncertain{unicité} ?}
\note{le premier ajout est cerclé dans la marge gauche, en face de
« Lemme » ; le second, en marge de la ligne suivante}

Ceci \ill{} ramène à prouver que \ill{} \ill{}. \struck{Soit} $U$
\struck{les \ill{}}. On peut supposer $S$ \uncertain{noethérien}. Il
suffit de prouver \uncertain{alors} \ill{} que l'on ait un schéma abélien
$A_U$ sur $U$, et $t$ un pt maximal de \ill{}, \ill{} un voisinage $V$ de
$t$. Si $U = T$, on \ill{} \ill{} voisinage ouvert $V$ de $t$,
\struck{\ill{}} $V$ un voisinage \ill{} un schéma abélien sur $V$, tel
que $A(s)$ provienne d'un schéma \ill{},
\struck{et tel que $V$ ne contienne \ill{} $x \in \mathrm{Ass}\,\mathcal{O}_S$
tel que $\overline{\{x\}} \not\ni t$, et tel que $B|U \cap V \simeq A_U$
\ill{} $s \in V$} [et \uncertain{tel} que $B \times_S
\mathrm{Spec}\,\mathcal{O}_{S,s} \simeq A(s)$ \uncertain{isom.}\
compatible avec les \ill{} $\,= U(s)$]. \struck{Soit} \ill{}
$s \in U \cap \mathrm{Spec}\,\mathcal{O}_{S,t}$, \ill{}
$x \in \mathrm{Ass}\,\mathcal{O}_S$, $\overline{x} \not\ni \{t\}
\Rightarrow x \notin V$, \uncertain{Car} \uncertain{supposons} que
$V \cap T \subset \overline{\{t\}}$. Soit alors $s \in V \cap T$, donc
$s \in \overline{\{t\}}$
\note{la phrase s'interrompt au bas du feuillet ; plusieurs lignes du
milieu sont biffées d'un trait et recouvertes d'additions interlinéaires,
dont l'ordre de lecture est incertain}

\page{18}
\note{numéroté par lui « 2 bis », cerclé, en haut à gauche ; papier bistre}

d'où $S(t) \to S(s)$.)

(\uncertain{On va} définir un isom.\ $A(s) \simeq B \times_S S(s)$, tel
que par changement de base $S(t) \to S(s)$, on obtienne l'isom.\
$A(t) \simeq B \times_S S(t)$ donné. \struck{\ill{}}

\textbf{Lemme.} $S$ schéma \uncertain{noeth.}\ [\uncertain{réduit}],
$t \in S$, $B$ schéma abélien sur $S$. Alors \struck{\ill{}} il existe un
voisinage ouvert $V$ de $t$, \struck{\ill{} $\forall s \in V$}, tel que
pour tout voisinage ouvert $V'$ de $t$, $V' \subset V$, tout schéma
abélien $B'$ sur $V'$, et tout isom.\
$\varphi_t : B' \times_S \mathrm{Spec}\,\mathcal{O}_{S,t} \simeq B
\times_S \mathrm{Spec}\,\mathcal{O}_{S,t}$, \uncertain{il y a} un isom.\
$\varphi(s)$ \ill{} \uncertain{voisinage} \ill{} \uncertain{unique} \ill{}
un isom.
\[
  \varphi : B' \simeq B|V'.
\]

L'\uncertain{unicité} sans \uncertain{généralité}, dès que toute
composante \uncertain{irréd.}\ de $V$ \uncertain{rencontre} $t$, i.e.\
que $V$ \ill{} \uncertain{les} comp.\ irréd.\ de $S$
\uncertain{contenant} \ill{} \uncertain{contenant} $t$.

\struck{Soit $S' \to S$} On peut supposer $S$ \uncertain{affine}.
(\uncertain{Soit} $S' \to S$ fini, tel que $S'$ réduit et géom.\
unibr.\ [\uncertain{supposons} d'abord $S$ réduit.]) \ill{}
$S' \times_S S(s) \to S(s)$ \uncertain{sont} \uncertain{birationnels}, et
\ill{}

Alors
\note{le feuillet s'arrête sur ce mot ; le reste est blanc}

\page{20}
\note{sans numéro de sa main ; une reprise du lemme du feuillet 3}

$A$, $B$ schémas abéliens sur $S$,
$u_p : T_p(A) \to T_p(B)$.
\note{en tête, à gauche, un petit croquis : $A$ et $B$ reliés par deux
traits à $S$ au-dessous}

\textbf{Question} : quand $u_p$ provient-il d'un hom $u : A \to B$.

\struck{(i) À cause de l'unicité, la descente fid.\ plate montre}

\struck{(ii) l'ens.\ $T$ des pts $s \in S$ tels que
$u_{ps} : T_p(A_s) \to T_p(B_s)$ provienne de $u_s : A_s \to B_s$
(condition géom.\ par (i)!) est \emph{fermé}.
\struck{Si $t \in T$ est tel que car.\ $k(t) =$}
(iii) Si $t \in T$ est tel que car $k(t) = p$, alors $\exists$ voisinage
ouvert $U$ de $t$ tel que $u_p|U$ soit algébrique, en particulier
\struck{$T$} $U$ est \uncertain{contenu} dans $T$.}
\note{(ii) et (iii) sont encadrés et biffés de quatre traits obliques}

a) Il y a un plus grand sous-schéma $T$ de $S$ sur lequel $u_p$ est
algébrique ; $T$ est fermé, sa formation commute à tout chgt de base
(nouvelles bases \uncertain{par} \uncertain{voisinages}
\uncertain{noethériens}).
\note{au-dessus de « plus grand », en interligne, une parenthèse d'une
lecture incertaine, \uncertain{(cf.\ Serre)}}

\note{sous le cadre, à gauche, un grand croquis : une région de la base
$S$ balayée de courbes parallèles, avec, marqués sur elle, « $S$ »,
« $S_0$ » (fortement repassé), « $S_1 \leftarrow S'$ » et, au-dessus,
« $A_{S'} \to B_{S'}$ » ; un point cerclé au coin inférieur gauche, où
aboutit une flèche}

b) Si $t \in T$ est tel que car $k(t) = p$, alors $T$ \uncertain{majore}
un voisinage ouvert de $t$. \struck{Donc si $S_p = V(p \cdot 1_S)$, alors
$T$ est un voisinage (\uncertain{schématique}) de $T \cap S_p$.}
\struck{Cor.\ Supposons que pour tout $s \in \mathrm{Ass}\,\mathcal{O}_S$,}
\struck{Supposons que pour tout $s \in$ \ill{} des pts de $S$ \ill{}}

\textbf{Cor.} Soit $Z$ l'ens.\ \struck{des \ill{}} \struck{(\ill{} pts
maximaux de $S$)} et de l'ens.\ des pts associés à $\mathcal{O}_S$
\emph{et} de \uncertain{ses} \uncertain{points} maximaux, \ill{} \ill{}
$\neq p$. Supposons que
\[
  z \in Z \Longrightarrow \exists\, s \in \overline{\{z\}} \text{ tel que }
  k(s) \text{ de car.\ } p \text{ et } s \in T
\]
(i.e.\ $\overline{\{z\}} \cap S_p \cap T \neq \emptyset$). Alors $u_p$ est
algébrique.
\note{la parenthèse « i.e.\ \ldots » est ajoutée au-dessus de la ligne ;
le « $\neq p$ » isolé en tête de ligne termine une incise biffée dont
l'attache est incertaine}

\end{document}
